\answer{ex:21:1}{(a)~$T=\tau_{\text{rec}}\ln\frac{1}{1-x^*}$：$1.84$~s和 $3.68$~s。(b)~从 $3.36$ 到 $4.24$~s：灵敏度 $\dd T/\dd x^*=\tau_{\text{rec}}/(1-x^*)=80$~s，$x^*\to1$ 时群放电不规则。}
\answer{ex:21:2}{(a)~$8.6$~W，与~\eqref{eq:energy} 一致。(b)~$205$~Mbit/s，$0.2$~W，是培养物（$10^{-4}$~W）的 $2\cdot10^3$ 倍。}
\answer{ex:21:3}{$x_{\text{ss}}=1/(1+Ur_b\tau_{\text{rec}})<x^*$ 的条件是 $r_b>(1/x^*-1)/(U\tau_{\text{rec}})=0.28$~Hz（$x^*=0.9$），$x^*=0.99$ 时为 $0.025$~Hz；突触强度降到 $x_{\text{ss}}$，即减弱 $10\%$ 和 $1\%$。}
\answer{ex:21:4}{$u(t-k)$ 彼此正交归一，$m_k=\|P_Vu(t-k)\|^2$，其中 $P_V$ 是到输出张成的 $N$ 维子空间上的投影；由贝塞尔不等式，$\sum_k m_k\le N$。}
\answer{ex:21:5}{有 $\tanh$ 时 $\mathrm{MC}\approx9$--$11$，线性储备池为 $15$--$18$（三个随机种子），都 $\le30$：非线性以记忆为代价。}
\answer{ex:21:6}{(a)~$n\le102$。(b)~高出 $5.6$ 个标准差。}
\answer{ex:21:7}{开放问题。关键对照：冻结读出，检查在无刺激的停顿之后改进是否依然存在；若不存在，则改变的是状态，而不是权重。}
